\documentclass[reqno]{amsart}
\usepackage[T1]{fontenc}
\usepackage{amsmath,amssymb,amsthm}
\usepackage{xr-hyper}
\usepackage[hidelinks]{hyperref}
\externaldocument[a1-]{20261126T080100--free-levels-and-short-bases__workbook}
\externaldocument[sd-]{20261125T000000--slope-method__workbook}
\externaldocument[kl-]{20261125T009000--kronecker-limit-and-chowla-selberg__workbook}
\externaldocument[tn-]{20261125T002100--definite-theta-lifts-and-norms__workbook}
\externaldocument[da-]{20261126T030300--the-data__workbook}
\externaldocument[ob-]{20261126T040100--spaces-forms-and-twists__workbook}
\externaldocument[ss-]{20261126T030400--the-seesaw-identity__workbook}
\externaldocument[dl-]{20261126T030500--divisibility-of-the-composite-lift__workbook}
\externaldocument[d2-]{20261126T050200--one-lift-at-a-small-non-split-prime__workbook}
\externaldocument[ta-]{20261126T030600--theorem-a-short-forms__workbook}

% Chapter number; \ref{chap} gives it.
\newcounter{chap}
% Chapter numbers, generated from 20261125T012000--grh-workbook__guide.org
% by scripts/make-chapters.py.  Do not edit: change the order of the
% index and rerun the script (./preview-site.sh and CI do this).
%
% Usage in a sheet, after \newcounter{chap}:
%   % Chapter numbers, generated from 20261125T012000--grh-workbook__guide.org
% by scripts/make-chapters.py.  Do not edit: change the order of the
% index and rerun the script (./preview-site.sh and CI do this).
%
% Usage in a sheet, after \newcounter{chap}:
%   % Chapter numbers, generated from 20261125T012000--grh-workbook__guide.org
% by scripts/make-chapters.py.  Do not edit: change the order of the
% index and rerun the script (./preview-site.sh and CI do this).
%
% Usage in a sheet, after \newcounter{chap}:
%   \input{chapters}\chapnumber{<sheet ID>}
\makeatletter
\providecommand\chapters@set[2]{\@namedef{chapters@#1}{#2}}
\chapters@set{20261125T001000}{1}% zeros-to-ramp-sums
\chapters@set{20261125T001100}{2}% inert-primes-and-estimate-b
\chapters@set{20261125T005000}{3}% galois-fourier-and-averaging
\chapters@set{20261125T005500}{4}% columns-from-a-sphere
\chapters@set{20261125T000000}{5}% slope-method
\chapters@set{20261126T030100}{6}% a-twist-for-each-prime
\chapters@set{20261125T003000}{7}% the-chain-and-the-square-root
\chapters@set{20261125T010000}{8}% rank-lemma-and-zero-estimate
\chapters@set{20261125T010500}{9}% poincare-series-and-near-orthogonality
\chapters@set{20261124T223000}{10}% elliptic-curves-for-the-proof
\chapters@set{20261125T009500}{11}% fock-model-and-theta-intertwiners
\chapters@set{20261124T223100}{12}% complex-multiplication-and-inert-reduction
\chapters@set{20261125T008000}{13}% one-cm-determinant
\chapters@set{20261124T214500}{14}% supersingular-theta-divisibility
\chapters@set{20261124T214600}{15}% the-isotropic-plane
\chapters@set{20261125T013700}{16}% pairings-on-a-formal-group-of-height-two
\chapters@set{20261125T013500}{17}% hecke-nondegeneracy-on-the-tree
\chapters@set{20261126T030200}{18}% the-bound-is-attained
\chapters@set{20261125T009000}{19}% kronecker-limit-and-chowla-selberg
\chapters@set{20261125T002000}{20}% definite-theta-lifts
\chapters@set{20261125T015500}{21}% siegel-weil-in-the-smallest-case
\chapters@set{20261125T002100}{22}% definite-theta-lifts-and-norms
\chapters@set{20261125T016500}{23}% hecke-operators-and-theta-lifts
\chapters@set{20261125T016700}{24}% divisibility-along-the-chain
\chapters@set{20261125T004000}{25}% denominators-and-sha
\chapters@set{20261125T004100}{26}% tate-twists-and-sha
\chapters@set{20261126T020430}{27}% the-minus-part-by-gauss-sums
\chapters@set{20261125T014100}{28}% congruence-modules
\chapters@set{20261125T014500}{29}% eisenstein-congruences-and-selmer-classes
\chapters@set{20261125T014700}{30}% pseudocharacters-and-the-length-bound
\chapters@set{20261125T011000}{31}% galois-representations-black-box
\chapters@set{20261126T030300}{32}% the-data
\chapters@set{20261126T020210}{33}% bessel-poincare-series-and-their-gram-operators
\chapters@set{20261126T030400}{34}% the-seesaw-identity
\chapters@set{20261126T030500}{35}% divisibility-of-the-composite-lift
\chapters@set{20261126T030600}{36}% theorem-a-short-forms
\chapters@set{20261126T030700}{37}% the-determinant-of-a-chain
\chapters@set{20261126T030800}{38}% the-cost-of-one-lift
\chapters@set{20261126T030900}{39}% denominators-of-projections-to-lifts
\chapters@set{20261126T031000}{40}% twists-with-few-congruences
\chapters@set{20261126T031100}{41}% theorem-c-and-the-main-inequality
\providecommand\chapnumber[1]{%
  \@ifundefined{chapters@#1}%
    {\PackageError{chapters}{Sheet #1 is not listed in the index}%
      {List it in 20261125T012000--grh-workbook__guide.org\MessageBreak
       and run scripts/make-chapters.py.}}%
    {\edef\thechap{\@nameuse{chapters@#1}}}}
\makeatother
\chapnumber{<sheet ID>}
\makeatletter
\providecommand\chapters@set[2]{\@namedef{chapters@#1}{#2}}
\chapters@set{20261125T001000}{1}% zeros-to-ramp-sums
\chapters@set{20261125T001100}{2}% inert-primes-and-estimate-b
\chapters@set{20261125T005000}{3}% galois-fourier-and-averaging
\chapters@set{20261125T005500}{4}% columns-from-a-sphere
\chapters@set{20261125T000000}{5}% slope-method
\chapters@set{20261126T030100}{6}% a-twist-for-each-prime
\chapters@set{20261125T003000}{7}% the-chain-and-the-square-root
\chapters@set{20261125T010000}{8}% rank-lemma-and-zero-estimate
\chapters@set{20261125T010500}{9}% poincare-series-and-near-orthogonality
\chapters@set{20261124T223000}{10}% elliptic-curves-for-the-proof
\chapters@set{20261125T009500}{11}% fock-model-and-theta-intertwiners
\chapters@set{20261124T223100}{12}% complex-multiplication-and-inert-reduction
\chapters@set{20261125T008000}{13}% one-cm-determinant
\chapters@set{20261124T214500}{14}% supersingular-theta-divisibility
\chapters@set{20261124T214600}{15}% the-isotropic-plane
\chapters@set{20261125T013700}{16}% pairings-on-a-formal-group-of-height-two
\chapters@set{20261125T013500}{17}% hecke-nondegeneracy-on-the-tree
\chapters@set{20261126T030200}{18}% the-bound-is-attained
\chapters@set{20261125T009000}{19}% kronecker-limit-and-chowla-selberg
\chapters@set{20261125T002000}{20}% definite-theta-lifts
\chapters@set{20261125T015500}{21}% siegel-weil-in-the-smallest-case
\chapters@set{20261125T002100}{22}% definite-theta-lifts-and-norms
\chapters@set{20261125T016500}{23}% hecke-operators-and-theta-lifts
\chapters@set{20261125T016700}{24}% divisibility-along-the-chain
\chapters@set{20261125T004000}{25}% denominators-and-sha
\chapters@set{20261125T004100}{26}% tate-twists-and-sha
\chapters@set{20261126T020430}{27}% the-minus-part-by-gauss-sums
\chapters@set{20261125T014100}{28}% congruence-modules
\chapters@set{20261125T014500}{29}% eisenstein-congruences-and-selmer-classes
\chapters@set{20261125T014700}{30}% pseudocharacters-and-the-length-bound
\chapters@set{20261125T011000}{31}% galois-representations-black-box
\chapters@set{20261126T030300}{32}% the-data
\chapters@set{20261126T020210}{33}% bessel-poincare-series-and-their-gram-operators
\chapters@set{20261126T030400}{34}% the-seesaw-identity
\chapters@set{20261126T030500}{35}% divisibility-of-the-composite-lift
\chapters@set{20261126T030600}{36}% theorem-a-short-forms
\chapters@set{20261126T030700}{37}% the-determinant-of-a-chain
\chapters@set{20261126T030800}{38}% the-cost-of-one-lift
\chapters@set{20261126T030900}{39}% denominators-of-projections-to-lifts
\chapters@set{20261126T031000}{40}% twists-with-few-congruences
\chapters@set{20261126T031100}{41}% theorem-c-and-the-main-inequality
\providecommand\chapnumber[1]{%
  \@ifundefined{chapters@#1}%
    {\PackageError{chapters}{Sheet #1 is not listed in the index}%
      {List it in 20261125T012000--grh-workbook__guide.org\MessageBreak
       and run scripts/make-chapters.py.}}%
    {\edef\thechap{\@nameuse{chapters@#1}}}}
\makeatother
\chapnumber{<sheet ID>}
\makeatletter
\providecommand\chapters@set[2]{\@namedef{chapters@#1}{#2}}
\chapters@set{20261125T001000}{1}% zeros-to-ramp-sums
\chapters@set{20261125T001100}{2}% inert-primes-and-estimate-b
\chapters@set{20261125T005000}{3}% galois-fourier-and-averaging
\chapters@set{20261125T005500}{4}% columns-from-a-sphere
\chapters@set{20261125T000000}{5}% slope-method
\chapters@set{20261126T030100}{6}% a-twist-for-each-prime
\chapters@set{20261125T003000}{7}% the-chain-and-the-square-root
\chapters@set{20261125T010000}{8}% rank-lemma-and-zero-estimate
\chapters@set{20261125T010500}{9}% poincare-series-and-near-orthogonality
\chapters@set{20261124T223000}{10}% elliptic-curves-for-the-proof
\chapters@set{20261125T009500}{11}% fock-model-and-theta-intertwiners
\chapters@set{20261124T223100}{12}% complex-multiplication-and-inert-reduction
\chapters@set{20261125T008000}{13}% one-cm-determinant
\chapters@set{20261124T214500}{14}% supersingular-theta-divisibility
\chapters@set{20261124T214600}{15}% the-isotropic-plane
\chapters@set{20261125T013700}{16}% pairings-on-a-formal-group-of-height-two
\chapters@set{20261125T013500}{17}% hecke-nondegeneracy-on-the-tree
\chapters@set{20261126T030200}{18}% the-bound-is-attained
\chapters@set{20261125T009000}{19}% kronecker-limit-and-chowla-selberg
\chapters@set{20261125T002000}{20}% definite-theta-lifts
\chapters@set{20261125T015500}{21}% siegel-weil-in-the-smallest-case
\chapters@set{20261125T002100}{22}% definite-theta-lifts-and-norms
\chapters@set{20261125T016500}{23}% hecke-operators-and-theta-lifts
\chapters@set{20261125T016700}{24}% divisibility-along-the-chain
\chapters@set{20261125T004000}{25}% denominators-and-sha
\chapters@set{20261125T004100}{26}% tate-twists-and-sha
\chapters@set{20261126T020430}{27}% the-minus-part-by-gauss-sums
\chapters@set{20261125T014100}{28}% congruence-modules
\chapters@set{20261125T014500}{29}% eisenstein-congruences-and-selmer-classes
\chapters@set{20261125T014700}{30}% pseudocharacters-and-the-length-bound
\chapters@set{20261125T011000}{31}% galois-representations-black-box
\chapters@set{20261126T030300}{32}% the-data
\chapters@set{20261126T020210}{33}% bessel-poincare-series-and-their-gram-operators
\chapters@set{20261126T030400}{34}% the-seesaw-identity
\chapters@set{20261126T030500}{35}% divisibility-of-the-composite-lift
\chapters@set{20261126T030600}{36}% theorem-a-short-forms
\chapters@set{20261126T030700}{37}% the-determinant-of-a-chain
\chapters@set{20261126T030800}{38}% the-cost-of-one-lift
\chapters@set{20261126T030900}{39}% denominators-of-projections-to-lifts
\chapters@set{20261126T031000}{40}% twists-with-few-congruences
\chapters@set{20261126T031100}{41}% theorem-c-and-the-main-inequality
\providecommand\chapnumber[1]{%
  \@ifundefined{chapters@#1}%
    {\PackageError{chapters}{Sheet #1 is not listed in the index}%
      {List it in 20261125T012000--grh-workbook__guide.org\MessageBreak
       and run scripts/make-chapters.py.}}%
    {\edef\thechap{\@nameuse{chapters@#1}}}}
\makeatother
\chapnumber{20261126T050300}
\renewcommand{\thesection}{\thechap.\arabic{section}}

\theoremstyle{definition}
\newtheorem{problem}{Problem}[section]
\newtheorem{definition}[problem]{Definition}
\theoremstyle{plain}
\newtheorem{theorem}[problem]{Theorem}
\theoremstyle{remark}
\newtheorem*{leanremark}{In the Lean formalization}

\newcommand{\Z}{\mathbb{Z}}
\newcommand{\R}{\mathbb{R}}
\newcommand{\C}{\mathbb{C}}
\newcommand{\gain}{\mathrm{gain}}
\newcommand{\rem}{\mathrm{rem}}
\newcommand{\hm}{h_{\mathrm{main}}}
\newcommand{\sinf}{s_\infty}
\newcommand{\sfin}{s_{\mathrm{fin}}}
\newcommand{\sfp}{s_{\mathrm{fin},p}}
\newcommand{\smain}{s_{\mathrm{main}}}
\newcommand{\Cch}{C_{\mathrm{ch}}}
\newcommand{\rd}{\mathrm{rd}_E}
\newcommand{\row}{\mathrm{row}}
\newcommand{\tr}{\mathfrak t}
\newcommand{\Aut}{\operatorname{Aut}}
\newcommand{\Sym}{\operatorname{Sym}}
\newcommand{\diag}{\operatorname{diag}}
\newcommand{\tl}[1]{\texttt{#1}}
\newcommand{\Oc}{\mathcal{O}}
\newcommand{\rs}{\mathrm{rs}}

\setlength{\emergencystretch}{3em}

\title{The size of the composite lift}
\author{}
\date{November 26, 2026}

\begin{document}
\maketitle
\refstepcounter{chap}\label{chap}

\paragraph{Prerequisites.}
Stirling's formula; exterior powers and the singular value decomposition of a complex matrix \cite[\S2.6]{HJ}.  From Chapter~\ref{da-chap}: the data (Definition~\ref{da-def:data}), the weights \(\kappa\), the prime sums \(\sfin\), the sizes \(\sinf\), \(\Omega\), \(\hm\), \(\smain\), \(\rem\) and \(U\) (Section~\ref{da-sec:amounts}), and the level primes (Section~\ref{da-sec:level}).  From Chapter~\ref{ob-chap}: the places \(\Phi\), \(\Phi_a\), the spaces \(X_j\), \(Y_j\), the source forms \(S\), whose values lie in Weyl modules with their Weyl (Gram) norms, the lift weights \(\kappa_j\), the gap \(\Delta = 2r^8\), and the diagram \(\beta(\tau, h)\) (Definition~\ref{ob-def:diagram}).  From Chapter~\ref{dl-chap}, Section~\ref{dl-sec:gain}: \(c_p\), \(I_p\), \(\sfp\) and \(\gain(p)\); and weak chain data (Definition~\ref{dl-def:weak}), the restricted chain \(R_{\vartheta,\Psi}\) (Definition~\ref{dl-def:restricted}) and the expansion of a chain datum into a weak one (Theorem~\ref{dl-thm:plain}).  From Chapter~\ref{ss-chap}, Section~\ref{ss-sec:theorem}: translations \(\tr\), translated vacua \(v_j(\tr)\), chain data and the composite lift \(B_{\vartheta,\Psi}\) with its restriction to \(Y_m\).  From Chapter~\ref{d2-chap}: the jet scale \(c_{\mathrm{jet}}\) (Definition~\ref{d2-def:jet-series}) and the kernel of one lift (Section~\ref{d2-sec:kernel}).  The Chowla--Selberg formula, Theorem~\ref{kl-thm:cs}; the Fock norm \(\lVert z^a\rVert^2 = a!\) of Problem~\ref{sd-pr:fock}; Chapter~\ref{tn-chap} for the smallest definite lifts.

\medskip
Theorem~A (Chapter~\ref{ta-chap}) compares the size of the composite lift with its divisibility.  Chapter~\ref{dl-chap} treats the divisibility; this sheet treats the size.  How large can \(B_{1,\Psi}\varphi\) be on every conjugate, if \(\varphi\) is small?  The answer has two halves.  At one pair of points the kernel of lift \(j\) multiplies conjugate Weyl norms by at most \(\exp(x^{50} + \rs_j/2)\), where \(\rs_j\) sums \(\sinf(l)\) over its rows at all places (Section~\ref{sec:kernel}): a Gaussian integral, and one factorial \(l!\) per row of length \(l\).  The finite prime sums cancel exactly against the gains (Problem~\ref{pr:exponent-identity}), and at each long place the level \(N\) pays for the excess of \(m\,\sinf(a)\) over \(m\,\smain\) up to \(32a\log b\) (Problem~\ref{pr:level-pays}).

\section{The exponent}\label{sec:exponent}

We recall the global prime sums and sizes from Chapter~\ref{da-chap} and the local gains from Chapter~\ref{dl-chap}.  The places \(\Phi\) are identified with \((\Z/q_1)^\times\), \([\nu]\) is the place of the residue class of an integer \(\nu\) if \(\gcd(\nu, q_1) = 1\), and the identity place otherwise, and \(\Phi_a\) is the set of \(a\)-places.  Recall \(\sinf(j) = \log j! + j\log(M/\Omega)\) with \(\Omega = e^{\mathrm{CS}(d)}\,2\pi/\rd\), \(c_1 = -1 - \mathrm{CS}(d)\), \(\smain = a(\log a - \frac1f\sum_{\tau \in \Phi}\log\kappa(\tau) + c_1)\), \(\hm = mk\sum_{\tau \in \Phi}(\sinf(\kappa(\tau)) - \sfin(\kappa(\tau), \tau))\) and \(\rem = \sum_{\tau \in \Phi_a}(\sfin(a, \tau) - \smain)\).  Let \(\Pi\) be the set of primes \(p < x_1\) with \(\chi_K(p) \ne 1\) (Section~\ref{dl-sec:gain}), and put, in Part~VIII,
\[
  \Sigma = \sum_{p \in \Pi}\gain(p)\log p .
\]
We use the following consequences of Chapter~\ref{da-chap} (Definition~\ref{da-def:data} and Problems~\ref{da-pr:rd}, \ref{da-pr:weights}, \ref{da-pr:jet} and~\ref{da-pr:MN}): \(x_1 \le a\); \(a/x_1 \le \kappa(\tau) \le b \le a\) for every \(\tau\); \(20\log b \le \log(M/\Omega) \le 23\log b\) (Problem~\ref{da-pr:omega}); \(\sinf(j) \ge 0\) for integers \(0 \le j \le a\); \(\log N \ge L_-(\kappa, m, M, \log b)\), as in Definition~\ref{da-def:data}; \(m(1 + \frac12\log a) \le a\log b\); and \(\log(M\rd/2\pi) \le 23\log b\).

By Problems~\ref{dl-pr:gain-zero} and~\ref{dl-pr:part-at-p}, and the sum over primes immediately after the latter,
\[
  2\Sigma = m\sum_{\tau \in \Phi_a}\sfin(a, \tau) + mk\sum_{\tau \in \Phi}\sfin(\kappa(\tau), \tau) .
\]

\begin{problem}\label{pr:exponent-identity}
\(\hm - m\,\rem + 2\Sigma = m\,\#\Phi_a\, \smain + mk\sum_{\tau \in \Phi}\sinf(\kappa(\tau))\).
\end{problem}

The prime sums have cancelled.

\begin{problem}\label{pr:stirling}
For every integer \(n \ge 1\), \(\log n! \le n\log n - n + 1 + \frac12\log n\).
\end{problem}

\begin{problem}\label{pr:excess}
\(\log(M/\Omega) = 1 + c_1 + \log(M\rd/2\pi)\), and
\[
  \sinf(a) - \smain \le a\Bigl(\frac1f\sum_{\tau \in \Phi}\log\kappa(\tau) + \log\frac{M\rd}{2\pi}\Bigr) + 1 + \tfrac12\log a .
\]
\end{problem}

\begin{problem}\label{pr:level-pays}
\(m(\sinf(a) - \smain) - a\log N \le 32\,a\log b\).
\end{problem}

Put \(\Cch = m\bigl(\#\Phi_a\,\sinf(a) + k\sum_{\tau \in \Phi}\sinf(\kappa(\tau))\bigr)\).

\begin{problem}\label{pr:chain-size}
\(\Cch - \#\Phi_a\,a\log N \le \hm - m\,\rem + 2\Sigma + 32\,fa\log b\).
\end{problem}

For \(n = 6\) the two sides of Problem~\ref{pr:stirling} compare as follows.

\begin{problem}\label{pr:stirling-six}
\(\log 6! + 6(1 - \log 6) = 1.8286\ldots \le 1 + \frac12\log 6 = 1.8958\ldots\).
\end{problem}

\section{One kernel in one coordinate}\label{sec:kernel}

Where does \(\sinf\) come from?  Moving a finite translate of the integral vacuum gives a finite translate of the moved vacuum times a scalar whose conjugates have modulus one.  For an automorphism over \(K\), the proof writes the conjugate of a theta jet of this latter translate as an integral of a conjugated polynomial against the Gaussian \(e^{-\sum_i\lvert y_i\rvert^2/2}\) and a bounded function \cite{GL}.  In this sheet \(l\) denotes the length of a row, \(L \ge 1\) a growth rate, and \(v\), \(t_0\), \(u\) are real numbers.

\begin{problem}\label{pr:factorial-device}
For \(\rho \ge 0\), \(l \in \mathbb{N}\) and \(v \in \R\), \(\rho^l \le l!\,\exp(lv + \rho e^{-v})\).
\end{problem}

\begin{problem}\label{pr:one-row}
In the one-coordinate model let \(l \ge 0\) be an integer, \(t_0 > 0\), \(L \ge 1\) and \(u \ge 0\).  If \(z \in \C\) and \(\lvert z\rvert^2 \le t_0(2L + 2)u\), then
\[
  \lvert z\rvert^l \le \exp\Bigl(\tfrac12\bigl(l/L + \log l! + l\log t_0\bigr) + Lu\Bigr) .
\]
\end{problem}

For a matrix the same device is applied to each row of a diagram.  In one variable the Fock norm \(\lVert z^l\rVert^2 = l!\) of Problem~\ref{sd-pr:fock} shows the same factorial.

\begin{problem}\label{pr:schur-model}
Let \(n \ge 1\) and let \(Z\) be a complex \(n \times n\) matrix and \(\rho_0 \ge \rho_1 \ge \dots \ge \rho_{n-1}\) the eigenvalues of \(Z^*Z\).  For \(1 \le h \le n\) the operator norm of \(\wedge^h Z\), for the inner product induced on \(\wedge^h\C^n\), satisfies \(\lVert\wedge^h Z\rVert^2 = \rho_0\rho_1\cdots\rho_{h-1}\).  For \(l \ge 1\), with the inner product on \(\Sym^l\C^n \subseteq (\C^n)^{\otimes l}\) induced from the tensor power, \(\lVert\Sym^l Z\rVert^2 \le \rho_0^l\).
\end{problem}

A column of height \(h\) carries \(\wedge^h\), a row of length \(l\) carries \(\Sym^l\); a general diagram mixes both, and the proof uses the following bound.  \emph{Schur norm.}  For each \(\tau\) in a finite set let \(z_\tau\) be a complex \(n'_\tau \times n_\tau\) matrix with \(n_\tau \le n'_\tau\), let \(b_{\tau,h}\) be numbers of columns of height \(h\), with \(b_{\tau,h} = 0\) for \(n_\tau < h < n'_\tau\), and put \(l_{\tau,i} = \sum_{i < h < n'_\tau} b_{\tau,h}\).  There are \(\rho_{\tau,i} \ge 0\), \(0 \le i < n_\tau\), with \(\sum_{i < n_\tau}\rho_{\tau,i}\) at most the sum of the squared absolute values of the entries of \(z_\tau\) for each \(\tau\), such that the Schur matrix of \((z_\tau)\) multiplies the squared Weyl norm of every vector by at most \(\prod_\tau\prod_{i < n_\tau}\rho_{\tau,i}^{l_{\tau,i}}\).  \emph{Status}: classical linear algebra (singular values); read from the proof \cite{GL}.

\begin{problem}\label{pr:gaussian}
Let \(L \ge 1\) and \(s' \in \R\).  Then \(Ls'/(2L + 2) - s'/2 = -s'/(2L + 2)\), and \(\frac1\pi\int_\C e^{-\lvert y\rvert^2/c}\,dA(y) = c\) for \(c > 0\).  For \(L = 2\), \(\int_0^\infty 2r e^{-r^2/6}\,dr = 6 = 2L + 2\).
\end{problem}

\begin{problem}\label{pr:integrated-row}
Let \(l \ge 0\) be an integer, \(t_0 > 0\), \(L \ge 1\), \(\mathrm{gr}(l) = \frac12(l/L + \log l! + l\log t_0)\), and \(H\colon \C \to \C\) measurable with \(\lvert H\rvert \le 1\).  Then
\[
  \Bigl\lvert\frac1\pi\int_\C(\sqrt{t_0}\,y)^l e^{-\lvert y\rvert^2/2}H(y)\,dA(y)\Bigr\rvert \le e^{\mathrm{gr}(l)}(2L + 2) ,
\]
and the Gaussian \(\exp(-\sum_i\lvert y_i\rvert^2/(2L + 2))\) in \(n_{\mathrm{c}}\) complex coordinates has integral \((2L + 2)^{n_{\mathrm{c}}}\) for the measure \(\pi^{-n_{\mathrm{c}}}dA\).
\end{problem}

The lift takes its Taylor series at the scale \(\pi_M c_{\mathrm{jet}}\), with \(c_{\mathrm{jet}}\) the jet scale of Definition~\ref{d2-def:jet-series}.  For each \(\sigma\) over \(K\) the conjugation law below gives a comparison scale \(t_\sigma\) with \(\lvert t_\sigma\rvert^2 = 1/\Omega\), and \(\lvert\sigma(\pi_M)\rvert^2 = M\); so the growth parameter is \(t_0 = \lvert\sigma(\pi_M)t_\sigma\rvert^2 = M/\Omega\), and \(2\,\mathrm{gr}(l) = l/L + \sinf(l)\).  Every automorphism of \(\C\) gives the same absolute values as one over \(K\) \cite{GL}.  With \(L = a + 1\), \(l/L \le 2\) for \(l \le 2a\).  The computation of \(\lvert t_\sigma\rvert^2\) uses the area of the conjugate lattice, the Chowla--Selberg unit, and the fact that the principal \(f\)-th root in Definition~\ref{d2-def:jet-series} has squared modulus \(\rd/\sqrt d\) on every conjugate over \(K\) \cite{GL}; the next two inputs enter it.

\emph{Conjugates of translated theta functions.}  Let \(d \equiv 7 \pmod 8\) be the prime of the data.  We call a lattice \(\Lambda \subset \C\) stable under \(\Oc_K\), with a point \(t' \in \C\) of exact order two modulo \(\Lambda\) such that \((1 \pm \sqrt{-d})t'/2 \notin \Lambda\), scaled so that \(\wp_\Lambda''(t') = 2\), with a chosen \(a_1\) such that \(a_1^2 = 12\wp_\Lambda(t')\), a \emph{normalized CM lattice}; for a lattice \(D\) write \(T^D_v f(z) = \exp(-\pi(z\bar v + \lvert v\rvert^2/2)/\operatorname{area}(D))f(z + v)\), and let \(\theta\) be the normalized theta section of \(\Lambda\) of \cite{GL} (the twisted translate by \(t'\) of the sigma function of \cite{GL}, which differs from the classical Weierstrass sigma function by a factor \(e^{cz^2}\), divided by the value of that translate at \(0\)).  Suppose \(\Lambda = \omega_0\Oc_K\), and let \(\sigma\) be an automorphism of \(\C\) over \(K\).  There are \(\omega' \in \C^\times\), a nonzero ideal \(\mathfrak a\) of \(\Oc_K\) and a finite adele \(\mathfrak b\) of \(K\) such that the lattice \(\Lambda'\) of the conjugate curve is \(\omega'\mathfrak a^2/\mathrm{Nm}\,\mathfrak a\), has area \(\lvert\omega'\rvert^2\sqrt d/2\) and contains \(n_0\omega'\) for some integer \(n_0 \ge 1\); and, with \(\theta'\) defined in the same way from the conjugate marked lattice, for every \(y \in K\) and every \(y' \in K\) close enough to \(\mathfrak b y\) in the finite adeles, applying \(\sigma\) coefficientwise to the Taylor series at \(0\) of \(T^\Lambda_{y\omega_0}\theta\) gives the Taylor series at \(0\) of \(T^{\Lambda'}_{y'\omega'}\theta'\).  \emph{Status}: this theta formulation is read from the proof \cite{GL}; for the classical main theorem of complex multiplication see \cite[Thm.~5.4]{Sh}.

\emph{The Chowla--Selberg unit.}  For every normalized CM lattice for \(K\), in any ideal class, there is \(u_{\mathrm{CS}} \in \C\) with \(\lvert\eta(u_{\mathrm{CS}})\rvert_p = 1\) for every prime \(p\) and every ring map \(\eta\colon \C \to \C_p\), such that for every automorphism \(\sigma\) of \(\C\) over \(K\) the area of the lattice of the conjugate curve times \(\lvert\sigma(u_{\mathrm{CS}})\rvert^2\) is \(\pi e^{\mathrm{CS}(d)}\,2\pi/\sqrt d\).  \emph{Status}: Theorem~\ref{kl-thm:cs} is the classical Chowla--Selberg product over ideal classes \cite{CS}; the statement above is read from the proof \cite{GL}.

Let \(\beta\) be the diagram of Definition~\ref{ob-def:diagram}.  Put \(a_\tau = a\) if \(\tau \in \Phi_a\) and \(a_\tau = 0\) otherwise.

\begin{definition}\label{def:rows}
Fix \(j < m\).  For \(\tau \in \Phi\) and \(0 \le i < n_j\) let
\[
  \row_\tau(i) = \sum_{i < h < n_{j+1}} \beta(\tau, h),
\]
the length of row \(i\), counted from \(0\), of the diagram at \(\tau\) formed by the columns of height less than \(n_{j+1}\).  We call these the \emph{row lengths of lift \(j\)}, drop \(j\) from the notation, and put \(\rs_j = \sum_{\tau \in \Phi}\sum_{0 \le i < n_j}\sinf(\row_\tau(i))\).
\end{definition}

Write \(G_f(V)\) for the group of finite adelic points of the unitary group of a space \(V\).  For \(j < m\), a translation \(\tr\), \(g \in G_f(X_j)\) and \(g' \in G_f(X_{j+1})\), consider the kernel of lift \(j\) of Section~\ref{d2-sec:kernel}, at the trivial twist and at the datum \(v_j(\tr)\) moved by \((g, g')\): a linear map from the Weyl vectors of the weight of \(X_j\) to those of \(X_{j+1}\), in which the values of forms lie.  The proof establishes \cite{GL}: \emph{on every conjugate this kernel multiplies Weyl norms by at most \(\exp(x^{50} + \rs_j/2)\)}; that is, for every value \(\omega\) and every \(\sigma \in \Aut(\C)\), the Weyl norm of \(\sigma\) applied to the image of \(\omega\) is at most \(\exp(x^{50} + \rs_j/2)\) times the Weyl norm of \(\sigma\omega\).  Here \(L = a + 1\), \(t_0 = M/\Omega\) and there are \(n_{\mathrm{c}} = fn_jn_{j+1}\) coordinates; besides \(\rs_j/2\) the exponent contains the logarithm of the vacuum bound, \(n_{\mathrm{c}}\log(2a + 4)\) and \(\sum_{\tau,i}\row_\tau(i)/(2L)\), and their sum is at most \(x^{50}\).

\section{The chain}\label{sec:chain}

The kernel bounds multiply along each tuple of classes, so their exponents involve \(\sum_{j<m}\rs_j\).  We bound this sum, count the classes, and restrict to \(Y_m\).

\begin{problem}\label{pr:jet-size-add}
For integers \(l, l_0, l_1 \ge 0\) with \(l \le l_0 + l_1\), \(\sinf(l) \le \sinf(l_0) + l_1\bigl(\log(l_0 + l_1 + 1) + \log(M/\Omega)\bigr)\).
\end{problem}

\begin{problem}\label{pr:row-bound}
Let \(j < m\), \(\tau \in \Phi\) and \(0 \le i < n_j\).  Then \(\row_\tau(0) \le a_\tau + b + r^{16}\) (the first-row slack in the proof is \(b + a/x_1 + r^{16}\)), \(\row_\tau(i) \le \kappa(\tau) + r^{16}\) for \(0 < i < k\), and \(\row_\tau(i) \le r^{16}\) for \(i \ge k\).  Each row length is at most \(2a\).
\end{problem}

\begin{problem}\label{pr:row-sum}
\(\sum_{j < m} \rs_j \le \Cch + U/2\).
\end{problem}

\emph{Class count.}  Theorem~\ref{a1-thm:class-count} bounds the number of classes at the standard level using short bases.  For the data, at a level containing the principal congruence subgroup of \(\pi_M\pi_N\), in rank at most \(n_m\), there are at most \(e^{x^{30}}\) classes.  \emph{Status}: read from the proof \cite{GL}.

\emph{Branching norm.}  The restriction from \(X_m\) to \(Y_m\) of Section~\ref{ss-sec:theorem} applies an integral branching matrix; its squared operator norm for the Weyl norms is at most \(e^{x^{70}}\).  \emph{Status}: read from the proof \cite{GL}.

For a form \(F\) on \(X_j\) or on \(Y_m\) and a ring homomorphism \(\sigma\colon \C \to \C\), \(\lvert\sigma F(g)\rvert^2\) denotes the Weyl squared norm of the conjugate value.

\begin{problem}\label{pr:restriction}
For a form \(F\) on \(X_m\), a ring homomorphism \(\sigma\colon \C \to \C\) and \(h \in G_f(Y_m)\), the restriction of \(F\) to \(Y_m\) satisfies
\[
  \lvert\sigma(\mathrm{res}\,F)(h)\rvert^2 \le \exp\Bigl(x^{70} - \log N\sum_{\tau \in \Phi}(a_\tau - mr^8)\Bigr)\,\lvert\sigma F(\diag(h, 1))\rvert^2 ,
\]
and \(\sum_{\tau \in \Phi}(a_\tau - mr^8) = \#\Phi_a\,a - fmr^8\).  Formally, for one place with \(a_\tau - mr^8 = 1\) and \(\pi_N = 1 + 2\sqrt{-1}\) (so \(N = 5\), which is not the level of any data), the scalar contributes \(\lvert 1/\bar\pi_N\rvert^2 = 1/5\) to the squared norm.
\end{problem}

\begin{problem}\label{pr:scale-check}
\(x^{70} + fmr^8\log N + 2(mx^{30} + x^{60} + mx^{50}) + 32fa\log b \le U/2\).
\end{problem}

Let \(\Psi\) be a chain datum for the trivial twist.  By Theorem~\ref{dl-thm:plain}, choose a weak chain datum \(\Psi'\) with \(B_{1,\Psi} = R_{1,\Psi'}\) on \(S\); then \(\sum_{\tr}\lvert\sigma(\Psi'(\tr))\rvert \le e^{x^{60}}\) for every \(\sigma \in \Aut(\C)\).  Fix \(\varphi \in S\) and \(y \in G_f(Y_m)\).  Write \(C_{\tr}\varphi\) for the chain of the \(m\) lifts of \(\varphi\) with the data \(v_j(\tr)\).  Before restriction the value at \(\diag(y, 1)\) is \(\sum_{\tr}\Psi'(\tr)(C_{\tr}\varphi)(\diag(y, 1))\), and the proof expands it as a sum over the \(\tr\) in the support of \(\Psi'\) and over the tuples of classes at the levels of \(X_0, \dots, X_{m-1}\), of \(\Psi'(\tr)\) times products of \(m\) kernels applied to values of \(\varphi\).  There are at most \(e^{mx^{30}}\) tuples of classes.  The resulting estimate for the restricted chain, whose exponent is controlled by Problems~\ref{pr:chain-size}, \ref{pr:row-sum}, \ref{pr:restriction} and~\ref{pr:scale-check}, gives the following theorem.  The problems prove the arithmetic of its exponent; the estimate for the restricted chain is read from the proof \cite{GL}.

\begin{theorem}\label{thm:size}
Let \(\Psi\) be a chain datum for the trivial twist, \(\sigma \in \Aut(\C)\), \(\varphi \in S\) and \(c \in \R\).  If \(\lvert\sigma\varphi(g)\rvert^2 \le e^c\) for every \(g \in G_f(X_0)\), then for every \(y \in G_f(Y_m)\)
\[
  \lvert\sigma(B_{1,\Psi}\varphi)(y)\rvert^2 \le \exp\bigl(c + \hm - m\,\rem + 2\Sigma + U\bigr) .
\]
\end{theorem}

The theorem is a statement at the trivial twist only.  The problems give the arithmetic estimates in its exponent; taken as given are the kernel bound with the Schur norm, the conjugation of translated theta functions and the Chowla--Selberg unit behind it, the class count at the congruence level, the branching norm, and the expansion of a chain datum.  In the exponent \(\Sigma\) has the coefficient \(2\); Chapter~\ref{ta-chap} sets the divisibility of Chapter~\ref{dl-chap} against it in Theorem~\ref{ta-thm:A}.

\begin{leanremark}
In the Lean proof \cite{GL} (prefix \tl{New.} omitted), Theorem~\ref{thm:size} is \tl{Data.weylNormSq\_compositeLift\_le}, deduced through \tl{Data.IsChainDatum.exists\_isWeakChainDatum\_compositeLift\_eq} from the bound \tl{Data.weylNormSq\_restrictedChain\_le} for the restricted chain, whose proof expands the chain by \tl{Data.chainCombination\_eq\_sum}.  The change in order of translation and the Weil action is made in \tl{Parameters.norm\_jet\_le\_of\_isVacuumCombination}.  For an automorphism over \(K\), the conjugate jet of the resulting translate of the moved vacuum is written as an integral in \tl{Parameters.map\_jet\_eq\_integral} (after \tl{Parameters.jet\_jetSeries\_mul}) and bounded in \tl{Parameters.norm\_jet\_translate\_le}, with the Taylor scale \(\pi_M c_{\mathrm{jet}}\) of \tl{Data.liftKernel}, the comparison scale from \tl{Parameters.exists\_isJetConjugation}, whose squared modulus \(1/\Omega\) is computed in \tl{Parameters.norm\_gal\_periodFactor\_div\_mul\_areaScale\_sq} inside \tl{Parameters.exists\_isJetConjugate\_translates}, the reduction to automorphisms over \(K\) in \tl{Parameters.exists\_algEquiv\_norm\_eq}, \(\lvert\sigma(\pi_M)\rvert^2 = M\) from \tl{Data.norm\_map\_star\_piMBar\_sq}, and the growth from \tl{sqrt\_prod\_pow\_le\_exp}; the first-row slack is in \tl{Data.diagramWeight\_le\_rowMainLength\_add\_rowSlack}; the restriction estimate is \tl{Data.weylNormSq\_restrictToY\_le}; the kernel bound of Section~\ref{sec:kernel} is \tl{Data.sqrt\_weylNormSq\_jetMap\_le}, with \(\rs_j\) the definition \tl{Data.rowSizeSum}.  The results taken as inputs on this sheet are \tl{CMLattice.complex\_multiplication\_theta} (with the sigma function \tl{PeriodPair.}\(\sigma\)), \tl{CMLattice.chowla\_selberg}, \tl{weylNormSq\_schur\_le}, \tl{Space.card\_stdClasses\_le}, applied with \(N' = q_1d\) in \tl{Data.card\_stdClasses\_le\_exp\_mul\_log}, with \tl{Data.card\_stdClasses\_le\_exp}, and \tl{exists\_branching\_matrix} with \tl{Data.isBranchingMap\_branching}.
\end{leanremark}

\section*{Hints}

\subsection*{\S\ref{sec:exponent}}
\begin{description}
\item[\ref{pr:exponent-identity}] The identity for \(2\Sigma\) displayed before Problem~\ref{pr:exponent-identity}.
\item[\ref{pr:stirling}] Show that \(\log n! - (n + \frac12)\log n + n\) decreases in \(n\); the series of \(\frac12\log\frac{1+t}{1-t}\) helps.
\item[\ref{pr:excess}] Problem~\ref{pr:stirling} with \(n = a\).
\item[\ref{pr:level-pays}] Compare Problem~\ref{pr:excess} with the lower bound for \(\log N\).  Which terms of size \(ma(\frac1f\sum\log\kappa + \log\frac{M\rd}{2\pi})\) cancel?
\item[\ref{pr:chain-size}] Problems~\ref{pr:exponent-identity} and~\ref{pr:level-pays}, and \(\#\Phi_a \le f\).
\end{description}

\subsection*{\S\ref{sec:kernel}}
\begin{description}
\item[\ref{pr:factorial-device}] One term of the exponential series.
\item[\ref{pr:one-row}] Apply Problem~\ref{pr:factorial-device} to \(\rho = \lvert z\rvert^2\) with \(v\) independent of \(z\); which \(v\) makes \((2L + 2)t_0e^{-v} \le 2L\)?
\item[\ref{pr:schur-model}] Singular value decomposition \(Z = P D Q^*\); \(\wedge^h\) and \(\Sym^l\) of a unitary matrix are unitary.
\item[\ref{pr:gaussian}] Polar coordinates.
\item[\ref{pr:integrated-row}] Problem~\ref{pr:one-row} with \(u = \lvert y\rvert^2/(2L + 2)\), then Problem~\ref{pr:gaussian}.
\end{description}

\subsection*{\S\ref{sec:chain}}
\begin{description}
\item[\ref{pr:jet-size-add}] Compare \(l!\) with \((l_0 + l_1)!\); what is the sign of \(\log(M/\Omega)\)?
\item[\ref{pr:row-bound}] Definition~\ref{ob-def:diagram}: which heights below \(n_{j+1}\) carry columns, and how many?  Use \(\kappa_1(\tau) = \kappa(\tau) + \frac32\Delta\), \(\kappa(\tau) \le b\) and the sizes of \(m\), \(r\), \(x\), \(b\) in Definition~\ref{da-def:data}.
\item[\ref{pr:row-sum}] Problem~\ref{pr:jet-size-add} with \(l_0 = a_\tau\) in row \(0\), \(l_0 = \kappa(\tau)\) for \(0 < i < k\) and \(l_0 = 0\) otherwise; how much can one extra box cost (Problem~\ref{da-pr:omega})?  Note \(k - 1 \le k\).  Then \(\sinf(0) = 0\), \(\sinf(\kappa) \ge 0\), and \(m \le 2\sqrt x\), \(a = xb\).
\item[\ref{pr:restriction}] The scalar of the restriction; \(\lvert\psi(\bar\pi_N)\rvert^2 = N\) for every embedding \(\psi\colon E \to \C\), since \(\pi_N\bar\pi_N = N\) and every complex embedding of \(E\) respects complex conjugation.  An integral matrix commutes with \(\sigma\).
\item[\ref{pr:scale-check}] Problem~\ref{da-pr:MN} for \(\log N\); compare each term with \(fa\log^3 b\), using \(b \ge x_1^{10^6}\).
\end{description}

\begin{thebibliography}{HJ}
\bibitem[Ba]{Ba} W.~Banaszczyk, \emph{New bounds in some transference theorems in the geometry of numbers}, Math.\ Ann.\ \textbf{296} (1993), 625--635.
\bibitem[CS]{CS} A.~Selberg and S.~Chowla, \emph{On Epstein's zeta-function}, J.\ Reine Angew.\ Math.\ \textbf{227} (1967), 86--110.
\bibitem[GL]{GL} L.~Sernau, \emph{GRH for Dirichlet \(L\)-functions, in Lean~4}, a shorter rewriting of the Lean proof, text by Claude, \url{https://github.com/lsernau/grh-lean}, commit \texttt{72fead1}, readable edition \texttt{grh/}.
\bibitem[HJ]{HJ} R.~A.~Horn and C.~R.~Johnson, \emph{Matrix Analysis}, 2nd ed., Cambridge University Press, 2013.
\bibitem[Sh]{Sh} G.~Shimura, \emph{Introduction to the Arithmetic Theory of Automorphic Functions}, Princeton University Press, 1971.
\end{thebibliography}

\end{document}
